\subsection{The inequality to be proved}
\label{subsec:central}
We now give the comparison in a form that separates the algebra from the numerical estimates. The notation introduced here will be used throughout the paper. Write
\begin{equation}\label{eq:normalization}
 W=\frac{R^2S}{3},\qquad
 W_1=2\sum_{(r,s)\in\Lambda}r,\qquad
 W_2=2\sum_{(r,s)\in\Lambda}s.
\end{equation}
The factor $W$ is a convenient normalization: the main logarithmic terms in both estimates have size $W\log D$. Direct summation gives
\begin{equation}\label{eq:radius-coefficient}
 \frac{W_1+W_2}{W}
 =1+\frac1K+\frac{3(K-1)}{2KS}+\frac{K-5}{2KS^2}.
\end{equation}

Choose $M$ columns $C$ for which the square submatrix $V_C$ is nonsingular, and put $\Delta_C=\det V_C$. The rank argument will supply such a choice; it will not specify the columns. The arithmetic estimates therefore have to hold for every possible choice $C$. They take the form
\begin{equation}\label{eq:arithmetic-interface}
 W\arith-2\sum_{(i,j)\in C} b(j)\ \leq\ \log|\Delta_C|^2.
\end{equation}
Here $\arith$ is a column-independent lower bound assembled from local divisibilities, and $b(j)$ is the loss from using a derivative of order $j$. Both depend on the chosen parameters. Section~\ref{sec:arithmetic} describes their construction and states the local estimate used below.

The main part of $b(j)$ is the first-prime-power sum
$\sum_q\lfloor j/q\rfloor\log q$. Legendre's formula shows that subtracting
it from $\log(j!)$ leaves higher-prime-power savings, with additional
contributions from non-inert primes. This remaining function is
nearly linear in $j$; we subtract a slope and bound its finite remainder. Set
\begingroup\postdisplaypenalty=10000
\begin{equation}\label{eq:F-H}
 F(j)=\log(j!)-b(j),\qquad H(j)=j\sigma-F(j),
\end{equation}
\endgroup
where $\sigma$ is a slope chosen in the arithmetic calculation. Let $c$ be any bound such that
\begin{equation}\label{eq:column-remainder}
 2\sum_{(i,j)\in C}H(j)\leq Mc
 \quad\hbox{for every }|C|=M.
\end{equation}
There are only $n$ columns of any one order $j$, so this is a finite selection problem with a capacity constraint. It need not be bounded by the much cruder $2M\max_jH(j)$.

Let $\rho$ contain all the nodes, put $u_i=z_i/\rho$, and define
\begin{equation}\label{eq:residual-matrix}
 B_{(r,s),(i,j)}(\ell)
 =\frac{\fall rj}{R^j}e^{\ell j}u_i^{r-j}\bar u_i^{s},
 \qquad \ell=\log(R/\rho)-\sigma.
\end{equation}
Here $\fall rj=r(r-1)\cdots(r-j+1)$ is the falling factorial, with
$\fall r0=1$; the entry is zero when $j>r$. The exact relation to the integral matrix is
\begin{equation}\label{eq:diagonal-scaling}
 B_C(\ell)
 =\operatorname{diag}_{\Lambda}(\rho^{-r-s})\,
 V_C\,
 \operatorname{diag}_{C}\bigl(j!(\rho/R)^je^{\ell j}\bigr).
\end{equation}
Taking squared absolute determinants, substituting \eqref{eq:F-H}, and then using \eqref{eq:arithmetic-interface} gives
\[
 W\arith\leq (W_1+W_2)\log\rho
       +\log|\det B_C(\ell)|^2+2\sum_C H(j).
\]
This calculation explains the normalization: the factorials are cancelled exactly, rather than estimated twice.

\begin{proposition}[Determinant comparison]\label{prop:comparison}
Suppose that $V$ has full row rank, that \eqref{eq:arithmetic-interface} and \eqref{eq:column-remainder} hold, and that
\[
 \frac1W\log|\det B_C(\ell)|^2\leq U_{\rm var}
 \qquad (|C|=M)
\]
holds uniformly over all column selections. Then
\begin{equation}\label{eq:master}
 \arith\leq \frac{W_1+W_2}{W}\log\rho
                  +U_{\rm var}+\frac{Mc}{W}.
\end{equation}
In particular, a strictly larger arithmetic lower bound is impossible.
\end{proposition}
\begin{proof}
Choose a nonsingular minor and apply the preceding identities. All bounds are uniform in the choice, so no information about which minor is nonsingular is needed.
\end{proof}

For the numerical comparisons we use the slightly weaker bound obtained
by adding $2M\log M/W$ to the right side of~\eqref{eq:master}.
This positive term is less than $3\times10^{-7}$ at the parameters for $499$.
Thus define
\begin{equation}\label{eq:metric-total}
 \metric=\frac{W_1+W_2}{W}\log\rho+U_{\rm var}
             +\frac{Mc+2M\log M}{W},
 \qquad G=\arith-\metric.
\end{equation}
A positive lower bound for $G$ rules out the proposed field. There are three separate tasks: prove full rank, establish the local arithmetic estimates, and bound the same minor at the complex embedding. The auxiliary-prime scalings used to prove rank in Section~\ref{sec:rank} do not enter \eqref{eq:diagonal-scaling}; after nonvanishing has been established, all size and divisibility estimates concern the original integral minor $\Delta_C$.

It is already possible to see why the number of nodes matters. The disk estimate gives $\rho\asymp n^{1/2}D^{1/4}$ once its boundary error is small. Reducing $n$ therefore lowers the common term in \eqref{eq:metric-total} by approximately
\[
 \frac12\left(1+\frac1K\right)\log\frac{n_{\rm old}}{n_{\rm new}}.
\]
At the two node counts in Table~\ref{tab:499}, this is about $0.085$. The remaining terms also change: there are fewer copies of each derivative order, and the proportion $\theta=M/(nJ)$ of available columns that must be selected increases. We call this proportion the \emph{occupancy}. Section~\ref{sec:metric} exploits these restrictions. In the actual calculation the combined gain is enough to change the sign of $G$.
